\section{Open and Responsible: los pesos abiertos son solo una cuarta parte}

Esta sección coloca la data governance en el centro del release de un modelo. La apertura comprende, como mínimo, datos accesibles, modelos accesibles, investigación reproducible y documentación completa; la responsibility exige además opt-out, PII, license, evaluation y limitation reporting.

\subsection{Cuatro niveles de apertura: Data, Models, Research, Documentation}

La slide divide la apertura en cuatro columnas. Los datos requieren inspection/opt-out/PII; los modelos requieren weights, training code y evaluation interface; la investigación requiere processing/training/evaluation reproducibility; la documentación requiere dataset/model cards, governance, architecture, cost y limitations.

\lecturefigure{slide-61.jpg}{La checklist de cuatro niveles de Open and Responsible Research on LLMs}{Deck oficial de Loubna Ben Allal, página 61}

\begin{importantbox}{Definición mínima de un release abierto}
\begin{enumerate}
  \item Poder saber qué data y qué transformation utilizó el modelo;
  \item Poder obtener, o al menos auditar, weights/code/config;
  \item Poder reproducir el training/evaluation path principal;
  \item Poder consultar license, limitations, PII/opt-out y cost.
\end{enumerate}
Contar solo con downloadable weights no basta para responder estas preguntas.
\end{importantbox}

\teachervoice{00:43:10--00:45:10, el ponente afirma explícitamente que un open release no consiste en «release model weights and stop there», sino en la combinación de data tools, opt-out, PII, reproducibility, evaluation tools y technical reports.}

\subsection{SantaCoder, StarCoder, StarCoder2: cómo itera el proyecto}

SantaCoder sirvió como ablation platform de alrededor de 1.1B; StarCoder escaló a 15B/alrededor de 1T tokens; StarCoder2 ofrece versiones de 3B/7B/15B y utiliza un corpus más grande y con más lenguajes. La slide muestra a la vez la disminución de la proporción de Python source y la continuidad de la transparencia y el acceso abierto.

\lecturefigure{slide-62.jpg}{De SantaCoder a StarCoder2: escala, lenguajes, proporción de Python y apertura}{Deck oficial de Loubna Ben Allal, página 62}

\begin{knowledgebox}{Por qué empezar con SantaCoder}
Una TB-scale data rule no puede probarse repetidamente de forma directa en un full run de 15B. El modelo pequeño ofrece una ablation platform para filter/format/objective; sin embargo, la extrapolation aún debe validarse en modelos más grandes, y no puede suponerse que todo data effect sea completamente scale-invariant.
\end{knowledgebox}

\subsection{Transparency Index: un indicador externo mide la divulgación, no toda la capacidad}

El Stanford Foundation Model Transparency Index asigna puntuaciones según dimensiones de divulgación como data, labor, compute, risks y downstream. En la instantánea presentada en clase, StarCoder aparece en los primeros lugares, lo que indica que el esfuerzo de documentación/transparencia fue captado por un marco externo; esto no demuestra que el modelo sea el más preciso ni el más seguro, ni que las controversias legales estén resueltas.

\lecturefigure{slide-63.jpg}{Resultados del Foundation Model Transparency Index al momento de la clase}{Deck oficial de Loubna Ben Allal, página 63}

\begin{warningbox}{Transparency y responsibility están relacionadas, pero no son equivalentes}
Divulgar malos resultados es más transparente que ocultarlos, pero aun así puede ser irresponsable; sin una divulgación completa puede haber controles internos, pero el público no puede verificarlos. El Index mide la información visible y no sustituye una independent audit ni un legal judgment.
\end{warningbox}

\subsection{Resumen de la sección}

Un responsible open model requiere publicar en conjunto data, weights, code, evaluation, documentation y governance. El valor de BigCode reside en convertir estos niveles en objetivos del proyecto; los indicadores de transparencia aportan una revisión externa, pero no sustituyen la evaluación de capacidades y riesgos.
